\heading{热学-24春-期末1}
\noteinfo{题目来源于赛艇，答案由AI生成。第 11 题缺失。}

\textbf{一、选择题}
\begin{question}
节流过程熵变化？
\begin{solution}
熵增，即 $\Delta S>0$。
\ansmath{\Delta S>0}
\end{solution}
\end{question}

\begin{question}
理想气体初态相同，分别经过可逆、不可逆绝热压缩到同一末压，比较末温。
\begin{solution}
可逆绝热满足 $\Delta S_R=0$，不可逆绝热满足 $\Delta S_{IR}>0$。对理想气体
\[
\Delta S=C_p\ln\frac{T_f}{T_i}-\nu R\ln\frac{p_f}{p_i}.
\]
在 $p_f/p_i$ 相同的条件下，$\Delta S_{IR}>\Delta S_R$ 推出
\[
T_{IR}>T_R.
\]
\ansmath{T_{IR}>T_R}
\end{solution}
\end{question}

\begin{question}
若气体满足 $\left(\partial U/\partial V\right)_T=-p$，则小过程是否有 $dS=0$？
\begin{solution}
由基本关系
\[
TdS=dU+p\,dV.
\]
若该小过程是等温过程，则
\[
dU=\left(\frac{\partial U}{\partial V}\right)_T dV=-p\,dV,
\]
于是 $TdS=0$，故 $dS=0$。
\ansmath{dU=\left(\frac{\partial U}{\partial V}\right)_T dV=-p\,dV,}
\end{solution}
\end{question}

\begin{question}
化学势 $\mu$ 的定义。
\begin{solution}
\[
\mu=\left(\frac{\partial G}{\partial N}\right)_{p,T}.
\]
\ansmath{\mu=\left(\frac{\partial G}{\partial N}\right)_{p,T}}
\end{solution}
\end{question}

\begin{question}
纯物质临界点附近的表面张力系数符号。
\begin{solution}
临界点附近表面张力趋于零，即 $\sigma\to 0$。
\ansmath{\sigma\to 0}
\end{solution}
\end{question}

\textbf{二、简答题}
\begin{question}
电解水反应 $\mathrm{H_2O(l)}\to \mathrm{H_2(g)}+\tfrac12\mathrm{O_2(g)}$，已知 298 K 下标准摩尔熵和 $\Delta_f H_m^\Theta$，求吸放热 $Q$ 与消耗电能 $E$。
\begin{solution}
先算熵变
\[
\Delta S_m^\Theta=S_m^\Theta(\mathrm{H_2})+\frac12S_m^\Theta(\mathrm{O_2})-S_m^\Theta(\mathrm{H_2O}).
\]
若过程在 $T=298\,\mathrm K$ 可逆进行，则
\[
Q=T\Delta S_m^\Theta.
\]
再由第一定律，电解所需总能量为焓变，故
\[
E=\Delta_f H_m^\Theta-Q=\Delta G_m^\Theta.
\]
\ansmath{E=\Delta_f H_m^\Theta-Q=\Delta G_m^\Theta}
\end{solution}
\end{question}

\begin{question}
证明等温线与绝热线不可能交于两点。
\begin{solution}
反证：若某条等温线与绝热线交于两点，则沿等温线和绝热线可围成一个循环。该循环要么从单一热源吸热并全部转化为功，要么违背 Clausius 不等式，分别与第一、第二定律矛盾，因此不可能。
\anstext{反证：若某条等温线与绝热线交于两点，则沿等温线和绝热线可围成一个循环。该循环要么从单一热源吸热并全部转化为功，要么违背 Clausius 不等式，分别与第一、第二定律矛盾，因此不可能。}
\end{solution}
\end{question}

\begin{question}
证明表面内能公式
\[
u_s=A_s\left(\sigma-T\left(\frac{\partial\sigma}{\partial T}\right)\right).
\]
\begin{solution}
表面自由能
\[
F_s=\sigma A_s.
\]
表面熵
\[
S_s=-\left(\frac{\partial F_s}{\partial T}\right)_{A_s}=-A_s\left(\frac{\partial\sigma}{\partial T}\right).
\]
故表面内能
\[
U_s=F_s+TS_s=A_s\left(\sigma-T\left(\frac{\partial\sigma}{\partial T}\right)\right).
\]
\ansmath{U_s=F_s+TS_s=A_s\left(\sigma-T\left(\frac{\partial\sigma}{\partial T}\right)\right)}
\end{solution}
\end{question}

\begin{question}
叙述一级相变特点及气液相平衡条件。
\begin{solution}
一级相变具有潜热和体积跃变；两相平衡条件为
\[
T_1=T_2,\qquad p_1=p_2,\qquad \mu_1=\mu_2.
\]
\ansmath{T_1=T_2,\qquad p_1=p_2,\qquad \mu_1=\mu_2}
\end{solution}
\end{question}

\begin{question}
叙述过热液体、过冷蒸气，并说明在等温线上的位置。
\begin{solution}
过热液体是已达到汽化条件却仍未汽化的液体，过冷蒸气是已达到液化条件却仍未液化的蒸气；它们都属亚稳态。在 vdW 等温线中，共存平台两侧的亚稳分支分别对应过热液体与过冷蒸气。
\anstext{过热液体是已达到汽化条件却仍未汽化的液体，过冷蒸气是已达到液化条件却仍未液化的蒸气；它们都属亚稳态。在 vdW 等温线中，共存平台两侧的亚稳分支分别对应过热液体与过冷蒸气。}
\end{solution}
\end{question}

\textbf{三、计算题}
\begin{question}
37 届决赛“量子热机”题。
\begin{solution}
上传稿未附题面，原稿仅注“自行搜索”，此处保留空位。
\anstext{题面缺失，答案从略。}
\end{solution}
\end{question}

\begin{question}
设离散体系满足 Maxwell--Boltzmann 分布
\[
p_i=\frac{e^{-E_i/(k_BT)}}{\sum_j e^{-E_j/(k_BT)}}.
\]
证明熵与自由能关系，并求熵函数。
\begin{solution}
配分函数 $Z=\sum_i e^{-E_i/(k_BT)}$，自由能
\[
F=-k_BT\ln Z.
\]
于是
\[
S=-\left(\frac{\partial F}{\partial T}\right)_V.
\]
再利用 $p_i=e^{-E_i/(k_BT)}/Z$，得
\[
F=U-TS,
\qquad U=\sum_i p_iE_i.
\]
化简即可得到
\[
S=-k_B\sum_i p_i\ln p_i.
\]
\ansmath{S=-k_B\sum_i p_i\ln p_i}
\end{solution}
\end{question}

\begin{question}
$1\,$mol vdW 气体节流过程，已知 $C_V,a,b,V_0,V_f,p_0,p_f$，求温度变化 $\Delta T$。
\begin{solution}
摩尔焓写为
\[
H_m=C_VT-\frac{a}{V_m}+pV_m+H_0.
\]
节流过程满足焓不变，故
\[
C_V(T_f-T_0)-\frac{a}{V_f}+p_fV_f+\frac{a}{V_0}-p_0V_0=0.
\]
因此
\[
\Delta T=T_f-T_0=\frac{a\left(\frac1{V_f}-\frac1{V_0}\right)-\left(p_fV_f-p_0V_0\right)}{C_V}.
\]
\ansmath{\Delta T=T_f-T_0=\frac{a\left(\frac1{V_f}-\frac1{V_0}\right)-\left(p_fV_f-p_0V_0\right)}{C_V}}
\end{solution}
\end{question}

\begin{question}
如图所示：左侧等容热容块热容 $C_r$，右侧理想气体有 $N$ 个粒子，热容记为 $C_g$；整体绝热，中间导热良好，左右同温。求过程方程及右侧从 $V_i$ 变到 $V_f$ 时的做功 $W$ 与传热 $Q_{rg}$。
\begin{solution}
右侧气体满足
\[
C_g dT+p\,dV=\delta Q_g,
\qquad p=\frac{Nk_BT}{V}.
\]
左侧热容块放热
\[
\delta Q_g=-dU_r=-C_r dT.
\]
故
\[
(C_r+C_g)dT+\frac{Nk_BT}{V}dV=0.
\]
积分得过程方程
\[
V\,T^{\frac{C_r+C_g}{Nk_B}}=\text{常数}
\qquad\text{或}\qquad
V^{Nk_B}T^{C_r+C_g}=\text{常数}.
\]
于是
\[
T_f=T_0\left(\frac{V_i}{V_f}\right)^{\frac{Nk_B}{C_r+C_g}}.
\]
做功可由
\[
dW=p\,dV=-(C_r+C_g)dT
\]
积分得
\[
W=(C_r+C_g)(T_0-T_f).
\]
若 $Q_{rg}$ 表示左侧传给右侧的热量，则
\[
Q_{rg}=C_r(T_0-T_f).
\]
\ansmath{Q_{rg}=C_r(T_0-T_f)}
\end{solution}
\end{question}

\begin{question}
三相（冰、水、水蒸气）各 $1\,$g，在 $t=0.01^\circ\mathrm C$ 下混合；给定热量 $Q$ 后求质量变化。
\begin{solution}
上传稿未给出完整数值与路径条件，无法唯一数值化。一般思路是：在三相点附近温度保持不变，所有热量先用于相变，按
\[
Q=L_f\,\Delta m_{\text{冰}\to\text{水}}+L_v\,\Delta m_{\text{水}\to\text{汽}}
\]
分段求解，并检查三相质量是否仍非负。
\ansmath{Q=L_f\,\Delta m_{\text{冰}\to\text{水}}+L_v\,\Delta m_{\text{水}\to\text{汽}}}
\end{solution}
\end{question}

\begin{question}
毛细管直径 $d$、高度 $h$，在高度 $h_1$ 处有侧支管口，完全浸润。求：
\begin{enumerate}[label=（\arabic*）,leftmargin=2.8em,itemsep=0.2em]
\item 毛细上升高度 $h$；
\item 侧口接触角；
\item 为什么不能构成永动机。
\end{enumerate}
\begin{solution}
\begin{enumerate}[label=（\arabic*）,leftmargin=2.8em,itemsep=0.2em]
\item 顶部液面压强
\[
p_t=p_0-\frac{4\sigma}{d}=p_0-\rho gh,
\]
故
\[
h=\frac{4\sigma}{\rho gd}.
\]
\item 侧口液面压强
\[
p_s=p_t+\rho g(h-h_1)=p_0-\rho gh_1<p_0.
\]
故液面向内凹，且
\[
\left|\Delta p_s\right|=\frac{4\sigma\cos\theta}{d}=\rho gh_1,
\]
因此
\[
\cos\theta=\frac{\rho gh_1d}{4\sigma}.
\]
\item 因侧口液面向内凹，水不会自发从侧口滴出，故不能持续输出重力势能，不构成永动机。
\end{enumerate}
\anstext{见上。}
\end{solution}
\end{question}

\begin{question}
已知压缩系数 $\beta$、表面张力 $\sigma$，求半径为 $r$ 的液滴相对平液面的密度比。
\begin{solution}
由 Laplace 压差
\[
\Delta p=\frac{2\sigma}{r}.
\]
压缩系数定义
\[
\beta=-\frac1V\frac{dV}{dp}.
\]
积分得
\[
\ln\frac{V_f}{V_i}=-\beta\Delta p
\quad\Rightarrow\quad
\ln\frac{\rho_i}{\rho_f}=-\beta\Delta p.
\]
当 $\beta\Delta p\ll1$ 时，
\[
\frac{\rho_f}{\rho_i}\approx 1+\beta\Delta p
=1+\frac{2\sigma\beta}{r}.
\]
\ansmath{\frac{\rho_f}{\rho_i}\approx 1+\beta\Delta p
=1+\frac{2\sigma\beta}{r}}
\end{solution}
\end{question}


